\documentclass[10pt,a4paper]{amsart}
\usepackage[margin=19mm]{geometry}
\usepackage{amsmath,amssymb,mathtools}
\usepackage[scr=rsfs,cal=euler]{mathalfa}
\usepackage{xfrac}
\usepackage[unicode,hidelinks,bookmarks=false]{hyperref}
\newcommand{\ie}{i.e.\ }
\newcommand{\cf}{cf.\ }
\newcommand{\resp}{respectively\ }
\newcommand{\Subsection}{\subsection}
\providecommand{\nobreakdash}{\nobreak\hbox{-}}
\setlength{\emergencystretch}{2em}
\multlinegap0pt
\usepackage{amssymb}
\usepackage{braket}
\usepackage{breqn}
\usepackage{dsfont}
\usepackage[shortlabels]{enumitem}
\usepackage{mathtools}
\usepackage{float}
\usepackage{xcolor}
\usepackage{tikz}
\newtheorem{definition}{Definition}
\newtheorem{assumption}{Assumption}
\newtheorem{proposition}{Proposition}
\newtheorem{lemma}{Lemma}
\def\N{\mathbb{N}}
\def\Om{\Omega}
\def\R{\mathbb{R}}
\def\argmin{\text{argmin}}
\def\de{\delta}
\def\ep{\epsilon}
\def\ex{\exists}
\def\grad{\nabla}
\def\liminf{\text{liminf}}
\def\vphi{\varphi}
\title{On convergence of solutions to nonlocal optimal control problems with quasi-minimization constraints}
\author{Joshua M. Siktar}
\date{Source: arXiv:2607.21822v1 (CC BY 4.0)}
\begin{document}
\maketitle

\noindent\emph{Source and licence.} On convergence of solutions to nonlocal optimal control problems with quasi-minimization constraints. Version arXiv:2607.21822v1 (CC BY 4.0), distributed by arXiv under the Creative Commons Attribution 4.0 International licence, http://creativecommons.org/licenses/by/4.0/, as recorded in the arXiv licence metadata for this exact version and shown on the abstract page https://arxiv.org/abs/2607.21822v1. Full licence notice: \texttt{licenses/arxiv-2607.21822-CC-BY-4.0.txt}.\par\smallskip

\noindent\textit{Editorial scope.} Counted unit: the article's lemma NLClosureOfMinimizationQuasi (source0.tex lines 635--637). The article's complete original proof of it is delivered with it (source0.tex lines 639--669, one \texttt{proof} environment of 31 lines). No mathematics is written, repaired or re-derived: the statement and the proof are the article's own lines, in the article's own notation, and no retained line is substituted or re-flowed.  Retained uncounted as the source-local context the proof needs: lines 468--473; lines 477--490; lines 502--504; lines 609--615.  Nothing was omitted from any retained span, so the package carries no omission record.  No retained line contains a citation, so the wrapper carries no bibliography and no bibliography entry.  No retained line contains a figure environment, an image-inclusion command or drawing code, so no image is delivered and the package declares no asset dependency.  Editorial formatting only, in addition: an AMS \texttt{amsart} preamble that restates the article's own declarations and macros used by the retained lines, namely the environments \texttt{definition}, \texttt{assumption}, \texttt{proposition}, \texttt{lemma} on separate counters, together with the standard packages those lines need.  The retained environments print in this document as Definition 1, Assumption 1, Proposition 1, Lemma 1 in order of appearance, whereas the article prints the same items with the numbering of its own full document.  The complete native source of the article as distributed in the arXiv e-print for this exact version is bundled unchanged as \texttt{sources/205895/source0.tex}.
    \begin{definition}[Quasiconvex energy]\label{Def: qConvexEnergy}
		An energy $W: \Om \times \R^n \times \R^{n \times n} \rightarrow \R$ is said to be \textbf{quasiconvex} if for a.e. $x \in \Om$, all $u \in \R^n$, and $A \in \R^{n \times n}$ the following inequality holds for all $\vphi \in W^{1, \infty}_0((0, 1)^n; \R^n)$:
		\begin{equation}\label{eq: qConvexEnergy}
			W(x, u, A) \ \leq \ \int_{(0, 1)^n}W(x, u, A + \grad \vphi(y))dy.
		\end{equation}
	\end{definition}

\begin{assumption}[Assumptions on energy density]
 \label{Assump: energyDensityAssump}
 The energy density $W: \Om \times \R^n \times \R^{n \times n} \to \R \cup \{ \infty \}$ is quasiconvex (in the sense of Definition \ref{Def: qConvexEnergy}) and is assumed to satisfy:
 \begin{itemize}
     \item $W$ is $\mathcal{L}^n \times \mathcal{B}^n \times \mathcal{B}^{n \times n}$-measurable and is a Carathéodory integrand %\footnote{$\mathcal{L}$ indicates Lebesgue measurability and $\mathcal{B}$ indicates Borel measurability.}.
\item $W (x, \cdot, \cdot)$ is lower semi-continuous for a.e.\ $x \in \Om$.
\item $W$ satisfies the growth bounds 
 \begin{equation} \label{Eq: WBound}
     c_w|A|^p - c_0 \ \leq \ W(x, z, A) \ \leq \ C_w(1 + |z|^p + |A|^p)
 \end{equation}
 for some $c_w, c_0, C_w > 0$, $p\in (1,\infty)$, a.e. $x \in \Om$, and all $A \in \R^{n \times n}$.% and some $b \in L^1(\Om)$ nonnegative. 
%\item$W(x, \cdot, \cdot)$ is of class $C^1$ for a.e. $x \in \Om$.
 \end{itemize}
 \end{assumption}

\begin{proposition}[Existence of minimizers for nonlocal energy density]\label{Prop: QuasiConvexityEnergy}
    Fix $g \in Z_{\text{ad}}$. Then there exists a $u_{\de, s} \in H^{s, p, \de}_0(\Om_{-\de}; \R^n)$ such that $u_{\de, s} \in \argmin_{v \in H^{s, p, \de}_0(\Om_{-\de}; \R^n)}\mathcal{W}^{\de, s}_g(v)$.
\end{proposition}

    \begin{lemma}[Boundedness of admissible states]\label{Lem: boundednessAdmissibleStatesNLQuasi}
        Let $W: \Om \times \R^n \times \R^{n \times n} \rightarrow \R$ be an energy density satisfying Assumption \ref{Assump: energyDensityAssump}, and let $\ep > 0$ be fixed. Then, the image of $\mathbb{T}^{\de, s}_{\ep}$ into $H^{s, p, \de}_0(\Om_{-\de}; \R^n)$, i.e., the collection of states which are $\ep$ quasi-minimizers of $\mathcal{W}^{\de, s}_g(\cdot)$ for some $g \in Z_{\text{ad}}$ given by
	\begin{equation}\label{Eq: UNLAdmisQuasi}
		V^{\de, s}_{\ep} \ := \ \{u_{\de, s} \in H^{s, p, \de}_0(\Om_{-\de}; \R^n) \ | \ \ex g_{\de, s} \in Z_{\text{ad}}, \ (u_{\de, s}, g_{\de, s}) \in \mathbb{T}^{\de, s}_{\ep}\},
	\end{equation}
	is a bounded subset of $H^{s, p, \de}_0(\Om_{-\de}; \R^n)$. 
    \end{lemma}

 \begin{lemma}[Weak closedness]\label{Lem: NLClosureOfMinimizationQuasi}
         Let $W: \Om \times \R^n \times \R^{n \times n} \rightarrow \R$ be an energy density satisfying Assumption \eqref{Assump: energyDensityAssump}. Then the set $\mathbb{T}^{\de, s}_{\ep}$ is weakly closed in $H^{s, p, \de}(\Om; \R^n) \times L^{p'}(\Om; \R^n)$.
    \end{lemma}

    \begin{proof}
    Let $\{(u_j, g_j)\}^{\infty}_{j = 1} \subset \mathbb{T}^{\de, s}_{\ep}$ be a sequence such that $u_j \rightharpoonup u$ weakly in $H^{s, p, \de}(\Om; \R^n)$ and $g_j \rightharpoonup g$ weakly in $L^{p'}(\Om; \R^n)$. We must show that $(u, g) \in \mathbb{T}^{\de, s}_{\ep}$, i.e.,
    \begin{equation}\label{Eq: NLClosureOfMinimizationQuasiEq1}
        \mathcal{W}^{\de, s}_g(u) - \min \mathcal{W}^{\de, s}_g(\cdot) \ \leq \ \ep.
    \end{equation}
    To this end, notice that for all $j \in \N^+$ we have
     \begin{equation}\label{Eq: NLClosureOfMinimizationQuasiEq2}
        \mathcal{W}^{\de, s}_{g_j}(u_j) - \min \mathcal{W}^{\de, s}_{g_j}(\cdot) \ \leq \ \ep.
    \end{equation}
    Due to the convergence properties of our sequences and weak lower semi-continuity of $W$, we have
    \begin{equation}\label{Eq: NLClosureOfMinimizationQuasiEq3}
        \liminf_{j \rightarrow \infty}\mathcal{W}^{\de, s}_{g_j}(u_j) \ \geq \  \mathcal{W}^{\de, s}_g(u).  
    \end{equation}
   Now we show that
 \begin{equation}\label{Eq: NLClosureOfMinimizationQuasiEq4}
        \lim_{j \rightarrow \infty}\min \mathcal{W}^{\de, s}_{g_j}(\cdot) \ = \ \min \mathcal{W}^{\de, s}_g(\cdot).
    \end{equation}
    Due to Proposition \ref{Prop: QuasiConvexityEnergy} applied for each $j \in \N^+$, there exists a sequence $\{v_j\}^{\infty}_{j = 1} \subset H^{s, p, \de}_0(\Om_{-\de}; \R^n)$ so that $\mathcal{W}^{\de,  s}_{g_j}(v_j) = \min\mathcal{W}^{\de, s}_{g_j}(\cdot)$. Then by Lemma \ref{Lem: boundednessAdmissibleStatesNLQuasi}, the family $\{v_j\}^{\infty}_{j = 1}$ is bounded in $H^{s, p, \de}_0(\Om_{-\de}; \R^n)$.  By reflexivity, there exists a (not relabeled) sub-sequence $\{v_j\}^{\infty}_{j = 1}$ and a limit $v \in H^{s, p, \de}_0(\Om_{-\de}; \R^n)$ so that $v_j \rightharpoonup v$ weakly in $H^{s, p, \de}(\Om; \R^n)$. Moreover, for each $j \in \N^+$ and $w \in H^{s, p, \de}_0(\Om_{-\de}; \R^n)$ arbitrary we have
    \begin{equation}\label{Eq: NLClosureOfMinimizationQuasiEq5}
        \mathcal{W}^{\de, s}_{g_j}(v_j) \ \leq \ \mathcal{W}^{\de, s}_{g_j}(w).
    \end{equation}
    Then by weak lower semi-continuity of $W$ and the weak convergence $g_j \rightharpoonup g$ in $L^{p'}(\Om; \R^n)$, we have
    \begin{equation}\label{Eq: NLClosureOfMinimizationQuasiEq6}
        \mathcal{W}^{\de,  s}_g(v) \ \leq \ \mathcal{W}^{\de, s}_g(w).
    \end{equation}    
    This means $\min \mathcal{W}^{\de, s}_g(\cdot) = \mathcal{W}^{\de, s}_g(v)$, and we conclude
    \begin{equation}\label{Eq: NLClosureOfMinimizationQuasiEq7}
       \min \mathcal{W}^{\de, s}_g(\cdot) \ \leq \ \liminf_{j \rightarrow \infty}\min \mathcal{W}^{\de, s}_{g_j}(\cdot).
    \end{equation}
To make \eqref{Eq: NLClosureOfMinimizationQuasiEq7} an equality, we may apply \eqref{Eq: NLClosureOfMinimizationQuasiEq5} by setting $w := v$ and sending $j \rightarrow \infty$, hence proving \eqref{Eq: NLClosureOfMinimizationQuasiEq4}. Finally, combining \eqref{Eq: NLClosureOfMinimizationQuasiEq2}, \eqref{Eq: NLClosureOfMinimizationQuasiEq3}, and \eqref{Eq: NLClosureOfMinimizationQuasiEq4} gives \eqref{Eq: NLClosureOfMinimizationQuasiEq1}, as desired.
    \end{proof}

\end{document}
